% GENERATED FILE. Edit canonical lessons, then run npm run generate:books.
% canonical-source-hash: fnv1a64:59aa3de60e4ce85d
% canonical-lessons: TE-C95-intelugu, TE-C95-howsay, TE-C95-gently, TE-C95-show, TE-C95-didnthear, TE-R95-second-pass-asking, TE-R95-second-pass-health

\chapter{How Do You Say It?}
\label{ch:te-how-do-you-say-it}
\hlchaptermodality{95}

\begin{chapteropening}
\textbf{By the end of this chapter:} \emph{I can ask how to say something in Telugu, ask to be shown, ask someone to speak gently, and say I could not hear.}

Complete TE-R95-second-pass-asking: courtesies and questions about the language, from cold.
\end{chapteropening}

\section[telugulō]{\te{తెలుగులో} (telugulō) — in Telugu}

\label{lesson:TE-C95-intelugu}



\begin{quote}
Before the new one: say the Telugu for need, then the Telugu for an introduction.
\end{quote}

\subsection*{You'll want to know: \te{తెలుగులో}}
\textbf{\te{తెలుగులో}} (\emph{telugulō}) — \textquotedblleft{}in Telugu\textquotedblright{}.

It is \textbf{\te{తెలుగు}} with \textbf{-\te{లో}}, \textquotedblleft{}in\textquotedblright{}, the ending you met on \textbf{\te{ఉండు}}. Any language takes it: \textbf{\te{ఇంగ్లీషులో}}, \textquotedblleft{}in English\textquotedblright{}.

\begin{grammarlens}[title={What you've built}]
A way to name the language you want an answer in.
\end{grammarlens}

\subsection*{Your turn}
\begin{itemize}
\raggedright
  \item \emph{Say it:} \emph{telugulō}
  \item \emph{Say it:} \emph{telugulō}, once more
  \item \emph{Say it:} say \emph{telugulō}, then \emph{iṅglīṣulō}
  \item \emph{Recall:} say the Telugu for need, then the Telugu for an introduction, then say \emph{telugulō} again
\end{itemize}

\begin{tcolorbox}[breakable,colback=teal!4,colframe=teal!35!black,title={Before you move on}]
What does \te{తెలుగులో} mean? (\textquotedblleft{}In Telugu\textquotedblright{}.) Say it once more.
\end{tcolorbox}
\section[elā aṇṭāru?]{\te{ఎలా} \te{అంటారు}? (elā aṇṭāru?) — how do you say it?}

\label{lesson:TE-C95-howsay}



\begin{quote}
Before the new one: say the Telugu for an introduction, then the Telugu for in Telugu.
\end{quote}

\subsection*{You'll want to know: \te{ఎలా} \te{అంటారు}?}
\textbf{\te{ఎలా} \te{అంటారు}?} (\emph{elā aṇṭāru?}) — \textquotedblleft{}how does one say it?\textquotedblright{}.

\textbf{\te{దీన్ని} \te{తెలుగులో} \te{ఎలా} \te{అంటారు}?} — \textquotedblleft{}how do you say this in Telugu?\textquotedblright{}. \textbf{\te{అను}}, \textquotedblleft{}to say\textquotedblright{}, is the verb inside \textbf{\te{అని}} and \textbf{\te{అనుకో}}.

\begin{grammarlens}[title={What you've built}]
The one question that lets a learner learn from anyone.
\end{grammarlens}

\subsection*{Your turn}
\begin{itemize}
\raggedright
  \item \emph{Say it:} \emph{elā aṇṭāru?}
  \item \emph{Say it:} \emph{elā aṇṭāru?}, once more
  \item \emph{Say it:} ask \emph{telugulō elā aṇṭāru?}
  \item \emph{Recall:} say the Telugu for an introduction, then the Telugu for in Telugu, then say \emph{elā aṇṭāru?} again
\end{itemize}

\begin{tcolorbox}[breakable,colback=teal!4,colframe=teal!35!black,title={Before you move on}]
What does \te{ఎలా} \te{అంటారు}? mean? (\textquotedblleft{}How do you say it?\textquotedblright{}.) Say it once more.
\end{tcolorbox}
\section[mellagā]{\te{మెల్లగా} (mellagā) — gently, softly, slowly}

\label{lesson:TE-C95-gently}



\begin{quote}
Before the new one: say the Telugu for in Telugu, then the Telugu for how do you say it.
\end{quote}

\subsection*{You'll want to know: \te{మెల్లగా}}
\textbf{\te{మెల్లగా}} (\emph{mellagā}) — \textquotedblleft{}gently, softly\textquotedblright{}, and \textquotedblleft{}slowly\textquotedblright{}.

It overlaps with \textbf{\te{నెమ్మదిగా}}. \textbf{\te{మెల్లగా}} leans to softness: a quiet voice, a careful hand.

\begin{grammarlens}[title={What you've built}]
A second slowly, with softness in it.
\end{grammarlens}

\subsection*{Your turn}
\begin{itemize}
\raggedright
  \item \emph{Say it:} \emph{mellagā}
  \item \emph{Say it:} \emph{mellagā}, once more
  \item \emph{Say it:} say \emph{mellagā māṭlāḍaṇḍi}
  \item \emph{Recall:} say the Telugu for in Telugu, then the Telugu for how do you say it, then say \emph{mellagā} again
\end{itemize}

\begin{tcolorbox}[breakable,colback=teal!4,colframe=teal!35!black,title={Before you move on}]
What does \te{మెల్లగా} mean? (\textquotedblleft{}Gently, softly, slowly\textquotedblright{}.) Say it once more.
\end{tcolorbox}
\section[cūpiñcaṇḍi]{\te{చూపించండి} (cūpiñcaṇḍi) — please show me}

\label{lesson:TE-C95-show}



\begin{quote}
Before the new one: say the Telugu for how do you say it, then the Telugu for gently.
\end{quote}

\subsection*{You'll want to know: \te{చూపించండి}}
\textbf{\te{చూపించండి}} (\emph{cūpiñcaṇḍi}) — \textquotedblleft{}please show me\textquotedblright{}.

It is built on \textbf{\te{చూడు}}, \textquotedblleft{}see\textquotedblright{}: to show is to \emph{make someone see}. When a word will not come, point and say \textbf{\te{చూపించండి}}.

\begin{grammarlens}[title={What you've built}]
A request for when words run out.
\end{grammarlens}

\subsection*{Your turn}
\begin{itemize}
\raggedright
  \item \emph{Say it:} \emph{cūpiñcaṇḍi}
  \item \emph{Say it:} \emph{cūpiñcaṇḍi}, once more
  \item \emph{Say it:} say \emph{cūpiñcaṇḍi}
  \item \emph{Recall:} say the Telugu for how do you say it, then the Telugu for gently, then say \emph{cūpiñcaṇḍi} again
\end{itemize}

\begin{tcolorbox}[breakable,colback=teal!4,colframe=teal!35!black,title={Before you move on}]
What does \te{చూపించండి} mean? (\textquotedblleft{}Please show me\textquotedblright{}.) Say it once more.
\end{tcolorbox}
\section[vinipiñcalēdu]{\te{వినిపించలేదు} (vinipiñcalēdu) — I could not hear it}

\label{lesson:TE-C95-didnthear}



\begin{quote}
Before the new one: say the Telugu for gently, then the Telugu for please show me.
\end{quote}

\subsection*{You'll want to know: \te{వినిపించలేదు}}
\textbf{\te{వినిపించలేదు}} (\emph{vinipiñcalēdu}) — \textquotedblleft{}I could not hear it\textquotedblright{}, literally \textquotedblleft{}it did not make itself heard\textquotedblright{}.

It ends in the \textbf{\te{లేదు}} you know. Pair it with a fix: \textbf{\te{వినిపించలేదు}, \te{గట్టిగా} \te{చెప్పండి}}.

\begin{grammarlens}[title={What you've built}]
In Telugu, how do you say, gently, show me, I could not hear: a learner's toolkit.
\end{grammarlens}

\subsection*{Your turn}
\begin{itemize}
\raggedright
  \item \emph{Say it:} \emph{vinipiñcalēdu}
  \item \emph{Say it:} \emph{vinipiñcalēdu}, once more
  \item \emph{Say it:} say \emph{vinipiñcalēdu, gaṭṭigā ceppaṇḍi}
  \item \emph{Recall:} say the Telugu for gently, then the Telugu for please show me, then say \emph{vinipiñcalēdu} again
\end{itemize}

\begin{tcolorbox}[breakable,colback=teal!4,colframe=teal!35!black,title={Before you move on}]
What does \te{వినిపించలేదు} mean? (\textquotedblleft{}I could not hear it\textquotedblright{}.) Say it once more.
\end{tcolorbox}
\section[(dialogue)]{Asking about the language — a second pass}

\label{lesson:TE-R95-second-pass-asking}



\begin{quote}
Point at something and ask how to say it in Telugu.
\end{quote}

\begin{grammarlens}[title={Grammar Lens: two kinds of polite}]
\noindent
\begin{tabularx}{\linewidth}{@{}>{\raggedright\arraybackslash}X>{\raggedright\arraybackslash}X@{}}
\toprule
\textbf{you ask} & \textbf{you answer} \\
\midrule
\textbf{\te{మన్నించండి}} & \textbf{\te{పరవాలేదు}} \\
\textbf{\te{టీ} \te{కావాలా}?} & \textbf{\te{అవసరం} \te{లేదు}} \\
\bottomrule
\end{tabularx}

An apology is met with \textquotedblleft{}no bother\textquotedblright{}; an offer is refused with \textquotedblleft{}no need\textquotedblright{}. Both answers end in \textbf{\te{లేదు}}.
\end{grammarlens}

\subsection*{Your turn}
Say these aloud:

\begin{itemize}
\raggedright
  \item beg pardon of an elder, then say you owe them an apology
  \item offer tea; then refuse help politely
  \item say you are glad to have met someone
  \item say you could not hear; ask them to speak gently, then to show you
\end{itemize}

\begin{tcolorbox}[breakable,colback=teal!4,colframe=teal!35!black,title={Before you move on}]
Pardon me? (\textbf{\te{మన్నించండి}}.) An apology? (\textbf{\te{క్షమాపణ}}.) Do you want it? (\textbf{\te{కావాలా}?}) No need? (\textbf{\te{అవసరం} \te{లేదు}}.) An introduction? (\textbf{\te{పరిచయం}}.) In Telugu? (\textbf{\te{తెలుగులో}}.) How do you say it? (\textbf{\te{ఎలా} \te{అంటారు}?}) Gently? (\textbf{\te{మెల్లగా}}.) Show me? (\textbf{\te{చూపించండి}}.) I could not hear? (\textbf{\te{వినిపించలేదు}}.)
\end{tcolorbox}
\section[(dialogue)]{From ill to well — a second pass}

\label{lesson:TE-R95-second-pass-health}



\begin{quote}
You are ill. Say you have a cold, you feel weak, and your body aches.
\end{quote}

\begin{grammarlens}[title={Grammar Lens: ill and well in one frame}]
\noindent
\begin{tabularx}{\linewidth}{@{}>{\raggedright\arraybackslash}X>{\raggedright\arraybackslash}X@{}}
\toprule
\textbf{not well} & \textbf{well again} \\
\midrule
\textbf{\te{నాకు} \te{నీరసంగా} \te{ఉంది}} & \textbf{\te{హాయిగా} \te{ఉంది}} \\
\textbf{\te{జలుబు} \te{చేసింది}} & \textbf{\te{క్షేమంగా} \te{ఉన్నారా}?} \\
\bottomrule
\end{tabularx}

The \textbf{-\te{గా} \te{ఉంది}} frame carries both halves: the weakness and the ease.
\end{grammarlens}

\subsection*{Your turn}
Say these aloud:

\begin{itemize}
\raggedright
  \item ask where the hospital is; say you have been hurt
  \item tell a friend to take the medicine
  \item a week later: ask an elder if they are keeping well, and a guest if they are comfortable
  \item agree gladly; then, when a word fails, ask to be shown, or say you could not hear
\end{itemize}

\begin{tcolorbox}[breakable,colback=teal!4,colframe=teal!35!black,title={Before you move on}]
A cold? (\textbf{\te{జలుబు}}.) Medicine? (\textbf{\te{మందు}}.) A hospital? (\textbf{\te{ఆసుపత్రి}}.) A wound? (\textbf{\te{గాయం}}.) Weakness? (\textbf{\te{నీరసం}}.) Wellbeing? (\textbf{\te{క్షేమం}}.) Comfortably? (\textbf{\te{హాయిగా}}.) Ease? (\textbf{\te{సుఖం}}.) Gladly? (\textbf{\te{సంతోషంగా}}.) The body? (\textbf{\te{ఒళ్ళు}}.)
\end{tcolorbox}
